\documentclass[../../main.tex]{subfiles}% 注意这里的文件路径不能用 ./main.tex ,否则用latexmk编译子文件会报错
\graphicspath{{\subfix{./image/}}{\subfix{../../image/}}} % 指定图片引用路径,后续可以直接使用图片文件名
% 如果图片存放在上级目录,则使用 ../../image/ 作为备用路径

% 例如:
% \begin{figure}[H]
% \centering
% \includegraphics[width=0.3\textwidth]{图.png}
% \caption{}
% \label{figure:图}
% \end{figure}
% 注意:上述\label{}一定要放在\caption{}之后,否则引用图片序号会只会显示??.

\begin{document}
% SOURCE p=79; printed=64; object=ch02-product-measures-title; region=upper
\section{乘积测度}\label{section:ch02-product-measures}
% SOURCE p=79; printed=64; object=ch02-product-measures-intro; region=upper
设 $(X,\mathscr{M},\mu)$ 和 $(Y,\mathscr{N},\nu)$ 为测度空间。我们已经在 $X\times Y$ 上定义了乘积 $\sigma$-代数 $\mathscr{M}\otimes\mathscr{N}$；现在要在其上构造一个测度，它在适当意义下是 $\mu$ 和 $\nu$ 的乘积。
% SOURCE p=79; printed=64; object=ch02-measurable-rectangle; region=upper
\begin{definition}[可测矩形]\label{definition:ch02-measurable-rectangle}
\textbf{可测矩形}是形如 $A\times B$ 的集合，其中 $A\in\mathscr{M}$，$B\in\mathscr{N}$。
\end{definition}
% SOURCE p=79; printed=64; object=ch02-rectangle-algebra; region=upper-middle
\begin{proposition}\label{proposition:ch02-rectangle-algebra}
有限个两两不交的可测矩形的并构成一个代数 $\mathscr{A}$，而 $\mathscr{A}$ 生成 $\mathscr{M}\otimes\mathscr{N}$。
\end{proposition}
\begin{proof}
由于
\begin{align*}
(A\times B)\cap(E\times F)&=(A\cap E)\times(B\cap F),\\
(A\times B)^c&=(X\times B^c)\cup(A^c\times B),
\end{align*}
结论由\cref{proposition:ch01-elementary-algebra}得出。
\end{proof}
% SOURCE p=79; printed=64; object=ch02-disjoint-rectangle-product; region=middle
\begin{proposition}\label{proposition:ch02-disjoint-rectangle-product}
若 $A\times B$ 是有限个或可数个两两不交的可测矩形 $A_j\times B_j$ 的并，则
\begin{align*}
\mu(A)\nu(B)=\sum_j\mu(A_j)\nu(B_j).
\end{align*}
\end{proposition}
\begin{proof}
有
\begin{align*}
\chi_A(x)\chi_B(y)=\chi_{A\times B}(x,y)
=\sum_j\chi_{A_j\times B_j}(x,y)=\sum_j\chi_{A_j}(x)\chi_{B_j}(y).
\end{align*}
关于 $x$ 积分并应用\cref{theorem:ch02-integral-series}，得到
\begin{align*}
\mu(A)\chi_B(y)
&=\int\chi_A(x)\chi_B(y)\,\mathrm{d}\mu(x)\\
&=\sum_j\int\chi_{A_j}(x)\chi_{B_j}(y)\,\mathrm{d}\mu(x)
=\sum_j\mu(A_j)\chi_{B_j}(y).
\end{align*}
再关于 $y$ 积分，即得所需等式。
\end{proof}
% SOURCE p=79; printed=64; object=ch02-product-premeasure; region=middle-lower
\begin{definition}[矩形代数上的预测度]\label{definition:ch02-product-premeasure}
若 $E\in\mathscr{A}$ 是两两不交的矩形 $A_1\times B_1,\ldots,A_n\times B_n$ 的并，定义
\begin{align*}
\pi(E)=\sum_{j=1}^n\mu(A_j)\nu(B_j),
\end{align*}
其中仍采用通常的约定 $0\cdot\infty=0$。
\end{definition}
% SOURCE p=79; printed=64; object=ch02-product-premeasure-well-defined; region=lower
\begin{proposition}\label{proposition:ch02-product-premeasure-well-defined}
$\pi$ 是良定义的，并且是 $\mathscr{A}$ 上的预测度。
\end{proposition}
\begin{proof}
$E$ 的任何两种有限矩形分解都有共同加细。因而，\cref{proposition:ch02-disjoint-rectangle-product}表明 $\pi$ 良定义，且为预测度。
\end{proof}
% SOURCE p=79; printed=64; object=ch02-product-measure; region=lower
\begin{definition}[乘积测度]\label{definition:ch02-product-measure}
由\cref{theorem:ch01-measure-extension}，$\pi$ 诱导一个外测度，其在 $\mathscr{M}\otimes\mathscr{N}$ 上的限制是延拓 $\pi$ 的测度。这个测度称为 $\mu$ 和 $\nu$ 的\textbf{乘积}，记作 $\mu\times\nu$。
\end{definition}
% SOURCE p=79; printed=64; object=ch02-product-sigma-finite-unique; region=bottom
\begin{proposition}\label{proposition:ch02-product-sigma-finite-unique}
若 $\mu$ 和 $\nu$ 为 $\sigma$-有限测度，则 $\mu\times\nu$ 也为 $\sigma$-有限测度，且是 $\mathscr{M}\otimes\mathscr{N}$ 上满足 $\mu\times\nu(A\times B)=\mu(A)\nu(B)$ 的唯一测度。
\end{proposition}
\begin{proof}
可写 $X=\bigcup_1^\infty A_j$、$Y=\bigcup_1^\infty B_k$，其中 $\mu(A_j)<\infty$、$\nu(B_k)<\infty$。于是 $X\times Y=\bigcup_{j,k}A_j\times B_k$，且 $\mu\times\nu(A_j\times B_k)<\infty$，故 $\mu\times\nu$ 为 $\sigma$-有限的。唯一性由\cref{theorem:ch01-measure-extension}得出。
\end{proof}
% SOURCE p=80; printed=65; object=ch02-finite-product-measure; region=top
\begin{definition}[有限个测度的乘积]\label{definition:ch02-finite-product-measure}
以上构造可以推广到任意有限个因子。若 $(X_j,\mathscr{M}_j,\mu_j)$ 为测度空间，$j=1,\ldots,n$，则将矩形定义为形如 $A_1\times\cdots\times A_n$ 的集合，其中 $A_j\in\mathscr{M}_j$。同样的方法构造出 $\mathscr{M}_1\otimes\cdots\otimes\mathscr{M}_n$ 上的\textbf{乘积测度} $\mu_1\times\cdots\times\mu_n$，它在这样的矩形上的值为 $\prod_1^n\mu_j(A_j)$。
\end{definition}
% SOURCE p=80; printed=65; object=ch02-product-associativity; region=upper
\begin{remark}
若各 $\mu_j$ 均为 $\sigma$-有限测度，则乘积测度由其在矩形上的取值唯一确定，而且乘积满足结合律。例如，将 $X_1\times X_2\times X_3$ 与 $(X_1\times X_2)\times X_3$ 等同。$\mathscr{M}_1\otimes\mathscr{M}_2\otimes\mathscr{M}_3$ 由矩形 $A_1\times A_2\times A_3$ 生成，而 $(\mathscr{M}_1\otimes\mathscr{M}_2)\otimes\mathscr{M}_3$ 由 $B\times A_3$ 生成，其中 $B\in\mathscr{M}_1\otimes\mathscr{M}_2$；这两个 $\sigma$-代数相同。此外，$\mu_1\times\mu_2\times\mu_3=(\mu_1\times\mu_2)\times\mu_3$，因为两者在三因子矩形上相同，故由唯一性处处相同。细节留作\cref{exercise:ch02-product-associativity}。以下结果都有显然的 $n$ 因子推广，但我们将只讨论 $n=2$ 的情形。
\end{remark}
% SOURCE p=80; printed=65; object=ch02-set-sections; region=middle
现在回到两个测度空间 $(X,\mathscr{M},\mu)$ 和 $(Y,\mathscr{N},\nu)$ 的情形。
\begin{definition}[集合的截面]\label{definition:ch02-set-sections}
若 $E\subseteq X\times Y$，对于 $x\in X$ 和 $y\in Y$，定义 $E$ 的 \textbf{$x$-截面} $E_x$ 和 \textbf{$y$-截面} $E^y$ 为
\begin{align*}
E_x=\{y\in Y:(x,y)\in E\},\qquad E^y=\{x\in X:(x,y)\in E\}.
\end{align*}
\end{definition}
% SOURCE p=80; printed=65; object=ch02-function-sections; region=middle
\begin{definition}[函数的截面]\label{definition:ch02-function-sections}
若 $f$ 是 $X\times Y$ 上的函数，则其 \textbf{$x$-截面} $f_x$ 和 \textbf{$y$-截面} $f^y$ 分别为 $Y$ 和 $X$ 上的函数，定义为 $f_x(y)=f^y(x)=f(x,y)$。
\end{definition}
% SOURCE p=80; printed=65; object=ch02-indicator-sections; region=middle
\begin{example}\label{example:ch02-indicator-sections}
$(\chi_E)_x=\chi_{E_x}$，$(\chi_E)^y=\chi_{E^y}$。
\end{example}
% SOURCE p=80; printed=65; object=ch02-section-measurability; number=2.34; region=middle-lower
\begin{proposition}\label{proposition:ch02-section-measurability}
\begin{enumerate}
\item\label{item:ch02-set-section-measurability} 若 $E\in\mathscr{M}\otimes\mathscr{N}$，则对所有 $x\in X$，$E_x\in\mathscr{N}$；对所有 $y\in Y$，$E^y\in\mathscr{M}$。
\item\label{item:ch02-function-section-measurability} 若 $f$ 为 $\mathscr{M}\otimes\mathscr{N}$-可测函数，则对所有 $x\in X$，$f_x$ 为 $\mathscr{N}$-可测函数；对所有 $y\in Y$，$f^y$ 为 $\mathscr{M}$-可测函数。
\end{enumerate}
\end{proposition}
% SOURCE p=80; printed=65; object=ch02-proof-section-measurability; region=lower
\begin{proof}
令 $\mathscr{R}$ 为所有满足 $E_x\in\mathscr{N}$、$E^y\in\mathscr{M}$（对一切 $x,y$）的 $E\subseteq X\times Y$ 所组成的族。$\mathscr{R}$ 包含所有可测矩形，因为例如 $(A\times B)_x$ 在 $x\in A$ 时为 $B$，在 $x\notin A$ 时为空集。又由于截面运算与并集、补集运算可交换，$\mathscr{R}$ 是一个 $\sigma$-代数，故包含 $\mathscr{M}\otimes\mathscr{N}$。这证明了\subcref{proposition:ch02-section-measurability}{item:ch02-set-section-measurability}。而\subcref{proposition:ch02-section-measurability}{item:ch02-function-section-measurability}随之成立，因为 $(f_x)^{-1}(B)=(f^{-1}(B))_x$，$f^y$ 的情形也相同。
\end{proof}
% SOURCE p=80; printed=65; object=ch02-monotone-class-intro; region=bottom
我们还需要一个技术性的引理。
% SOURCE p=80; printed=65; object=ch02-monotone-class; region=bottom
\begin{definition}[单调类]\label{definition:ch02-monotone-class}
集合 $X$ 上的\textbf{单调类}是一个族 $\mathscr{C}\subseteq\mathscr{P}(X)$，它对可数递增并和可数递减交封闭。也就是说，若 $E_j\in\mathscr{C}$ 且 $E_1\subseteq E_2\subseteq\cdots$，则 $\bigcup_1^\infty E_j\in\mathscr{C}$；对于递减序列的交也有相应结论。
\end{definition}
% SOURCE p=80,81; printed=65,66; object=ch02-generated-monotone-class; region=bottom/top
\begin{definition}[生成的单调类]\label{definition:ch02-generated-monotone-class}
显然，每个 $\sigma$-代数都是单调类，任意一族单调类的交仍是单调类。因此，对任意 $\mathscr{E}\subseteq\mathscr{P}(X)$，都存在唯一的包含 $\mathscr{E}$ 的最小单调类，称为由 $\mathscr{E}$ \textbf{生成}的单调类。
\end{definition}
% SOURCE p=81; printed=66; object=ch02-monotone-class-lemma; number=2.35; region=upper
\begin{lemma}[单调类引理]\label{lemma:ch02-monotone-class}
若 $\mathscr{A}$ 是 $X$ 的子集构成的代数，则由 $\mathscr{A}$ 生成的单调类 $\mathscr{C}$ 与由 $\mathscr{A}$ 生成的 $\sigma$-代数 $\mathscr{M}$ 相同。
\end{lemma}
% SOURCE p=81; printed=66; object=ch02-proof-monotone-class; region=upper-middle
\begin{proof}
由于 $\mathscr{M}$ 是单调类，有 $\mathscr{C}\subseteq\mathscr{M}$；若能证明 $\mathscr{C}$ 是 $\sigma$-代数，就有 $\mathscr{M}\subseteq\mathscr{C}$。为此，对 $E\in\mathscr{C}$ 定义
\begin{align*}
\mathscr{C}(E)=\{F\in\mathscr{C}:E\setminus F,\ F\setminus E\text{ 和 }E\cap F\text{ 均属于 }\mathscr{C}\}.
\end{align*}
显然，$\varnothing$ 和 $E$ 属于 $\mathscr{C}(E)$，且 $E\in\mathscr{C}(F)$ 当且仅当 $F\in\mathscr{C}(E)$。此外，容易验证 $\mathscr{C}(E)$ 是单调类。若 $E\in\mathscr{A}$，则对所有 $F\in\mathscr{A}$ 都有 $F\in\mathscr{C}(E)$，因为 $\mathscr{A}$ 是代数；即 $\mathscr{A}\subseteq\mathscr{C}(E)$，从而 $\mathscr{C}\subseteq\mathscr{C}(E)$。因此，若 $F\in\mathscr{C}$，则对所有 $E\in\mathscr{A}$，都有 $F\in\mathscr{C}(E)$。但这意味着对所有 $E\in\mathscr{A}$，都有 $E\in\mathscr{C}(F)$，故 $\mathscr{A}\subseteq\mathscr{C}(F)$，从而 $\mathscr{C}\subseteq\mathscr{C}(F)$。于是，若 $E,F\in\mathscr{C}$，则 $E\setminus F$ 和 $E\cap F$ 均属于 $\mathscr{C}$。由于 $X\in\mathscr{A}\subseteq\mathscr{C}$，$\mathscr{C}$ 因而是代数。然而，若 $\{E_j\}_1^\infty\subseteq\mathscr{C}$，则对所有 $n$，有 $\bigcup_1^n E_j\in\mathscr{C}$；又由于 $\mathscr{C}$ 对可数递增并封闭，可得 $\bigcup_1^\infty E_n\in\mathscr{C}$。所以 $\mathscr{C}$ 是 $\sigma$-代数，证明完成。
\end{proof}
% SOURCE p=81; printed=66; object=ch02-section-integrals-intro; region=middle
现在来讨论本节的主要结果，它们将 $X\times Y$ 上的积分与 $X$、$Y$ 上的积分联系起来。
% SOURCE p=81; printed=66; object=ch02-product-section-integral; number=2.36; region=middle
\begin{theorem}\label{theorem:ch02-product-section-integral}
设 $(X,\mathscr{M},\mu)$ 和 $(Y,\mathscr{N},\nu)$ 为 $\sigma$-有限测度空间。若 $E\in\mathscr{M}\otimes\mathscr{N}$，则函数 $x\mapsto\nu(E_x)$ 和 $y\mapsto\mu(E^y)$ 分别在 $X$ 和 $Y$ 上可测，且
\begin{align*}
\mu\times\nu(E)=\int\nu(E_x)\,\mathrm{d}\mu(x)=\int\mu(E^y)\,\mathrm{d}\nu(y).
\end{align*}
\end{theorem}
% SOURCE p=81,82; printed=66,67; object=ch02-proof-product-section-integral; region=lower/top
\begin{proof}
先假设 $\mu$ 和 $\nu$ 有限，令 $\mathscr{C}$ 为使定理结论成立的所有 $E\in\mathscr{M}\otimes\mathscr{N}$ 所组成的集合。若 $E=A\times B$，则 $\nu(E_x)=\chi_A(x)\nu(B)$，$\mu(E^y)=\mu(A)\chi_B(y)$，故显然 $E\in\mathscr{C}$。由可加性，有限个两两不交的矩形的并属于 $\mathscr{C}$；因而由\cref{lemma:ch02-monotone-class}，只需证明 $\mathscr{C}$ 是单调类。若 $\{E_n\}$ 为 $\mathscr{C}$ 中的递增序列，且 $E=\bigcup_1^\infty E_n$，则函数 $f_n(y)=\mu((E_n)^y)$ 可测，且逐点递增趋于 $f(y)=\mu(E^y)$。因此 $f$ 可测，由单调收敛定理，
\begin{align*}
\int\mu(E^y)\,\mathrm{d}\nu(y)
&=\lim\int\mu((E_n)^y)\,\mathrm{d}\nu(y)\\
&=\lim\mu\times\nu(E_n)=\mu\times\nu(E).
\end{align*}
同样，$\mu\times\nu(E)=\int\nu(E_x)\,\mathrm{d}\mu(x)$，故 $E\in\mathscr{C}$。类似地，若 $\{E_n\}$ 为 $\mathscr{C}$ 中的递减序列，且 $E=\bigcap_1^\infty E_n$，则函数 $y\mapsto\mu((E_1)^y)$ 属于 $L^1(\nu)$，因为 $\mu((E_1)^y)\le\mu(X)<\infty$ 且 $\nu(Y)<\infty$；于是可应用控制收敛定理证明 $E\in\mathscr{C}$。故 $\mathscr{C}$ 是单调类，有限测度空间的情形得证。

最后，若 $\mu$ 和 $\nu$ 为 $\sigma$-有限测度，可以将 $X\times Y$ 写成有限测度矩形的递增序列 $\{X_j\times Y_j\}$ 的并。若 $E\in\mathscr{M}\otimes\mathscr{N}$，对每个 $j$，把上述论证用于 $E\cap(X_j\times Y_j)$，得到
\begin{align*}
\mu\times\nu(E\cap(X_j\times Y_j))
&=\int\chi_{X_j}(x)\nu(E_x\cap Y_j)\,\mathrm{d}\mu(x)\\
&=\int\chi_{Y_j}(y)\mu(E^y\cap X_j)\,\mathrm{d}\nu(y).
\end{align*}
最后再应用一次单调收敛定理，就得到所需结论。
\end{proof}
% SOURCE p=82; printed=67; object=ch02-fubini-tonelli; number=2.37; region=upper-middle
\begin{theorem}[Fubini--Tonelli 定理]\label{theorem:ch02-fubini-tonelli}
设 $(X,\mathscr{M},\mu)$ 和 $(Y,\mathscr{N},\nu)$ 为 $\sigma$-有限测度空间。
\begin{enumerate}
\item\label{item:ch02-tonelli} （Tonelli）若 $f\in L^+(X\times Y)$，则函数 $g(x)=\int f_x\,\mathrm{d}\nu$ 和 $h(y)=\int f^y\,\mathrm{d}\mu$ 分别属于 $L^+(X)$ 和 $L^+(Y)$，且
% SOURCE p=82; printed=67; object=ch02-fubini-tonelli-equation; number=2.38; region=middle
\begin{align}\label{align:ch02-fubini-tonelli}
\int f\,\mathrm{d}(\mu\times\nu)
&=\int\left[\int f(x,y)\,\mathrm{d}\nu(y)\right]\mathrm{d}\mu(x)\notag\\
&=\int\left[\int f(x,y)\,\mathrm{d}\mu(x)\right]\mathrm{d}\nu(y).
\end{align}
% SOURCE p=82; printed=67; object=ch02-fubini-tonelli-part-b; region=middle
\item\label{item:ch02-fubini} （Fubini）若 $f\in L^1(\mu\times\nu)$，则对于几乎处处的 $x\in X$，$f_x\in L^1(\nu)$；对于几乎处处的 $y\in Y$，$f^y\in L^1(\mu)$。几乎处处有定义的函数 $g(x)=\int f_x\,\mathrm{d}\nu$ 和 $h(y)=\int f^y\,\mathrm{d}\mu$ 分别属于 $L^1(\mu)$ 和 $L^1(\nu)$，且\eqref{align:ch02-fubini-tonelli}成立。
\end{enumerate}
\end{theorem}
% SOURCE p=82; printed=67; object=ch02-proof-fubini-tonelli; region=middle-lower
\begin{proof}
当 $f$ 为示性函数时，Tonelli 定理归结为\cref{theorem:ch02-product-section-integral}；因此，由线性性，它对非负简单函数也成立。若 $f\in L^+(X\times Y)$，根据\cref{theorem:ch02-simple-approximation}，取逐点递增趋于 $f$ 的简单函数序列 $\{f_n\}$。单调收敛定理首先表明，相应的 $g_n$ 和 $h_n$ 递增趋于 $g$ 和 $h$（从而 $g$ 和 $h$ 可测）；其次表明
\begin{align*}
\int g\,\mathrm{d}\mu&=\lim\int g_n\,\mathrm{d}\mu=\lim\int f_n\,\mathrm{d}(\mu\times\nu)=\int f\,\mathrm{d}(\mu\times\nu),\\
\int h\,\mathrm{d}\nu&=\lim\int h_n\,\mathrm{d}\nu=\lim\int f_n\,\mathrm{d}(\mu\times\nu)=\int f\,\mathrm{d}(\mu\times\nu).
\end{align*}
这正是\eqref{align:ch02-fubini-tonelli}。这就证明了 Tonelli 定理，同时也说明：若 $f\in L^+(X\times Y)$ 且 $\int f\,\mathrm{d}(\mu\times\nu)<\infty$，则 $g<\infty$、$h<\infty$ 几乎处处成立，即对几乎处处的 $x$，$f_x\in L^1(\nu)$；对几乎处处的 $y$，$f^y\in L^1(\mu)$。若 $f\in L^1(\mu\times\nu)$，将这些结果用于 $f$ 的实部与虚部的正部和负部，便得到 Fubini 定理的结论。
\end{proof}
% SOURCE p=82,83; printed=67,68; object=ch02-fubini-tonelli-remarks; region=bottom/top
\Needspace{7\baselineskip}
\begin{remark}
这里作几点说明。
\begin{itemize}
\item 通常省略\eqref{align:ch02-fubini-tonelli}中累次积分的方括号，即
\begin{align*}
\int\left[\int f(x,y)\,\mathrm{d}\mu(x)\right]\mathrm{d}\nu(y)
=\iint f(x,y)\,\mathrm{d}\mu(x)\,\mathrm{d}\nu(y)=\iint f\,\mathrm{d}\mu\,\mathrm{d}\nu.
\end{align*}
\item $\sigma$-有限性假设是必要的，见\cref{exercise:ch02-fubini-non-sigma-finite}。
\item 假设 $f\in L^+(X\times Y)$ 或 $f\in L^1(\mu\times\nu)$ 在两个方面都是必要的。首先，即使 $f$ 不为 $\mathscr{M}\otimes\mathscr{N}$-可测函数，也可能对所有 $x,y$，$f_x$ 和 $f^y$ 都可测，且累次积分 $\iint f\,\mathrm{d}\mu\,\mathrm{d}\nu$ 和 $\iint f\,\mathrm{d}\nu\,\mathrm{d}\mu$ 都存在。但是，这两个累次积分未必相等，见\cref{exercise:ch02-nonmeasurable-separately-measurable}。其次，若 $f$ 不非负，则即使 $\int|f|\,\mathrm{d}(\mu\times\nu)=\infty$，也可能对所有 $x,y$，$f_x$ 和 $f^y$ 都可积，而且两个累次积分都存在。但它们仍未必相等，见\cref{exercise:ch02-nonintegrable-iterated}。
\item Fubini 定理和 Tonelli 定理经常配合使用。通常，我们希望交换二重积分 $\iint f\,\mathrm{d}\mu\,\mathrm{d}\nu$ 中的积分次序。首先利用 Tonelli 定理将 $\int|f|\,\mathrm{d}(\mu\times\nu)$ 算成累次积分，从而验证其小于 $\infty$；然后应用 Fubini 定理，得到 $\iint f\,\mathrm{d}\mu\,\mathrm{d}\nu=\iint f\,\mathrm{d}\nu\,\mathrm{d}\mu$。例子见\cref{section:ch02-n-dimensional-lebesgue-integral}中的习题。
\end{itemize}
\end{remark}
% SOURCE p=83; printed=68; object=ch02-product-incomplete; region=middle
\begin{remark}
即使 $\mu$ 和 $\nu$ 都完备，$\mu\times\nu$ 也几乎从不完备。事实上，假设存在非空集 $A\in\mathscr{M}$，使 $\mu(A)=0$，并且 $\mathscr{N}\ne\mathscr{P}(Y)$。（例如，$\mu=\nu$ 为 $\mathbb{R}$ 上的 Lebesgue 测度时就是如此。）若 $E\in\mathscr{P}(Y)\setminus\mathscr{N}$，则由\subcref{proposition:ch02-section-measurability}{item:ch02-set-section-measurability}，$A\times E\notin\mathscr{M}\otimes\mathscr{N}$；但 $A\times E\subseteq A\times Y$，且 $\mu\times\nu(A\times Y)=0$。
\end{remark}
% SOURCE p=83; printed=68; object=ch02-complete-product-intro; region=middle
如果希望使用完备测度，当然可以考虑 $\mu\times\nu$ 的完备化。在这种情形下，$X\times Y$ 上函数的可测性与其 $x$-截面、$y$-截面的可测性之间，关系不再如此简单。不过，将 Fubini--Tonelli 定理作适当改述后，它仍然成立。
% SOURCE p=83; printed=68; object=ch02-complete-fubini-tonelli; number=2.39; region=middle-lower
\begin{theorem}[完备测度的 Fubini--Tonelli 定理]\label{theorem:ch02-complete-fubini-tonelli}
设 $(X,\mathscr{M},\mu)$ 和 $(Y,\mathscr{N},\nu)$ 为完备的 $\sigma$-有限测度空间，$(X\times Y,\mathscr{L},\lambda)$ 为 $(X\times Y,\mathscr{M}\otimes\mathscr{N},\mu\times\nu)$ 的完备化。若 $f$ 为 $\mathscr{L}$-可测函数，并且满足下列条件之一：
\begin{enumerate}
\item\label{item:ch02-complete-tonelli} $f\ge0$；
\item\label{item:ch02-complete-fubini} $f\in L^1(\lambda)$，
\end{enumerate}
则对几乎处处的 $x$，$f_x$ 为 $\mathscr{N}$-可测函数；对几乎处处的 $y$，$f^y$ 为 $\mathscr{M}$-可测函数。在\subcref{theorem:ch02-complete-fubini-tonelli}{item:ch02-complete-fubini}的情形中，$f_x$ 和 $f^y$ 还分别对几乎处处的 $x$ 和 $y$ 可积。此外，$x\mapsto\int f_x\,\mathrm{d}\nu$ 和 $y\mapsto\int f^y\,\mathrm{d}\mu$ 可测，在\subcref{theorem:ch02-complete-fubini-tonelli}{item:ch02-complete-fubini}的情形中还可积，且
\begin{align*}
\int f\,\mathrm{d}\lambda=\iint f(x,y)\,\mathrm{d}\mu(x)\,\mathrm{d}\nu(y)
=\iint f(x,y)\,\mathrm{d}\nu(y)\,\mathrm{d}\mu(x).
\end{align*}
\end{theorem}
% SOURCE p=83; printed=68; object=ch02-complete-fubini-tonelli-proof-reference; region=lower
这个定理是\cref{theorem:ch02-fubini-tonelli}的一个相当容易的推论；证明提纲见\cref{exercise:ch02-complete-fubini-proof}。

\subsection*{习题}
% SOURCE p=83; printed=68; object=ch02-exercise-product-associativity; number=45; region=lower
\begin{exercise}\label{exercise:ch02-product-associativity}
若 $(X_j,\mathscr{M}_j)$ 为可测空间，$j=1,2,3$，则 $\bigotimes_{j=1}^3\mathscr{M}_j=(\mathscr{M}_1\otimes\mathscr{M}_2)\otimes\mathscr{M}_3$。此外，若 $\mu_j$ 为 $(X_j,\mathscr{M}_j)$ 上的 $\sigma$-有限测度，则 $\mu_1\times\mu_2\times\mu_3=(\mu_1\times\mu_2)\times\mu_3$。
\end{exercise}
% SOURCE p=83,84; printed=68,69; object=ch02-exercise-fubini-non-sigma-finite; number=46; region=bottom/top
\begin{exercise}\label{exercise:ch02-fubini-non-sigma-finite}
设 $X=Y=[0,1]$，$\mathscr{M}=\mathscr{N}=\mathscr{B}_{[0,1]}$，$\mu$ 为 Lebesgue 测度，$\nu$ 为计数测度。若 $D=\{(x,x):x\in[0,1]\}$ 为 $X\times Y$ 中的对角线，则 $\iint\chi_D\,\mathrm{d}\mu\,\mathrm{d}\nu$、$\iint\chi_D\,\mathrm{d}\nu\,\mathrm{d}\mu$ 和 $\int\chi_D\,\mathrm{d}(\mu\times\nu)$ 三者互不相等。（要计算 $\int\chi_D\,\mathrm{d}(\mu\times\nu)=\mu\times\nu(D)$，需回到 $\mu\times\nu$ 的定义。）
\end{exercise}
% SOURCE p=84; printed=69; object=ch02-exercise-nonmeasurable-separately-measurable; number=47; region=upper
\begin{exercise}\label{exercise:ch02-nonmeasurable-separately-measurable}
设 $X=Y$ 为不可数线性序集，且对每个 $x\in X$，$\{y\in X:y<x\}$ 都可数。
\begin{example*}
可数序数组成的集合。
\end{example*}
设 $\mathscr{M}=\mathscr{N}$ 为可数集或余可数集组成的 $\sigma$-代数，并在 $\mathscr{M}$ 上定义 $\mu=\nu$：当 $A$ 可数时，$\mu(A)=0$；当 $A$ 余可数时，$\mu(A)=1$。令 $E=\{(x,y)\in X\times X:y<x\}$。则对所有 $x,y$，$E_x$ 和 $E^y$ 都可测，而且 $\iint\chi_E\,\mathrm{d}\mu\,\mathrm{d}\nu$ 与 $\iint\chi_E\,\mathrm{d}\nu\,\mathrm{d}\mu$ 均存在但不相等。（若接受连续统假设，可取 $X=[0,1]$，赋予一个非标准序，从而得到集合 $E\subseteq[0,1]^2$，使对所有 $x,y$，$E_x$ 可数且 $E^y$ 余可数［特别地，均为 Borel 集］，但 $E$ 不 Lebesgue 可测。）
\end{exercise}
% SOURCE p=84; printed=69; object=ch02-exercise-nonintegrable-iterated; number=48; region=upper-middle
\begin{exercise}\label{exercise:ch02-nonintegrable-iterated}
设 $X=Y=\mathbb{N}$，$\mathscr{M}=\mathscr{N}=\mathscr{P}(\mathbb{N})$，$\mu=\nu$ 为计数测度。定义 $f(m,n)$：当 $m=n$ 时为 $1$，当 $m=n+1$ 时为 $-1$，其余情形为 $0$。则 $\int|f|\,\mathrm{d}(\mu\times\nu)=\infty$，且 $\iint f\,\mathrm{d}\mu\,\mathrm{d}\nu$ 与 $\iint f\,\mathrm{d}\nu\,\mathrm{d}\mu$ 均存在但不相等。
\end{exercise}
% SOURCE p=84; printed=69; object=ch02-exercise-complete-fubini-proof; number=49; region=middle
\begin{exercise}\label{exercise:ch02-complete-fubini-proof}
利用\cref{theorem:ch02-fubini-tonelli}、\cref{proposition:ch02-measurable-version-completion}及下列引理，证明\cref{theorem:ch02-complete-fubini-tonelli}。
\begin{enumerate}
\item\label{item:ch02-null-sections} 若 $E\in\mathscr{M}\times\mathscr{N}$ 且 $\mu\times\nu(E)=0$，则对几乎处处的 $x$ 和 $y$，有 $\nu(E_x)=\mu(E^y)=0$。
\item\label{item:ch02-null-function-sections} 若 $f$ 为 $\mathscr{L}$-可测函数，且 $f=0$ $\lambda$-几乎处处成立，则对几乎处处的 $x$ 和 $y$，$f_x$ 和 $f^y$ 可积，且 $\int f_x\,\mathrm{d}\nu=\int f^y\,\mathrm{d}\mu=0$。（这里需要 $\mu$ 和 $\nu$ 的完备性。）
\end{enumerate}
\end{exercise}
% SOURCE p=84; printed=69; object=ch02-exercise-area-under-graph; number=50; region=middle-lower
\begin{exercise}\label{exercise:ch02-area-under-graph}
设 $(X,\mathscr{M},\mu)$ 为 $\sigma$-有限测度空间，$f\in L^+(X)$。令
\begin{align*}
G_f=\{(x,y)\in X\times[0,\infty]:y\le f(x)\}.
\end{align*}
则 $G_f$ 为 $\mathscr{M}\times\mathscr{B}_{\mathbb{R}}$-可测集，且 $\mu\times m(G_f)=\int f\,\mathrm{d}\mu$；若将 $G_f$ 定义中的不等式 $y\le f(x)$ 替换为 $y<f(x)$，结论也相同。（为证明 $G_f$ 的可测性，注意映射 $(x,y)\mapsto f(x)-y$ 为 $(x,y)\mapsto(f(x),y)$ 和 $(z,y)\mapsto z-y$ 的复合。）这就是微积分中熟知的“函数的积分是其图像下的面积”这一定理的精确表述。
\end{exercise}
% SOURCE p=84; printed=69; object=ch02-exercise-product-functions; number=51; region=lower
\begin{exercise}\label{exercise:ch02-product-functions}
设 $(X,\mathscr{M},\mu)$ 和 $(Y,\mathscr{N},\nu)$ 为任意测度空间（不一定 $\sigma$-有限）。
\begin{enumerate}
\item\label{item:ch02-product-function-measurable} 若 $f:X\to\mathbb{C}$ 为 $\mathscr{M}$-可测函数，$g:Y\to\mathbb{C}$ 为 $\mathscr{N}$-可测函数，且 $h(x,y)=f(x)g(y)$，则 $h$ 为 $\mathscr{M}\otimes\mathscr{N}$-可测函数。
\item\label{item:ch02-product-function-integral} 若 $f\in L^1(\mu)$ 且 $g\in L^1(\nu)$，则 $h\in L^1(\mu\times\nu)$，且 $\int h\,\mathrm{d}(\mu\times\nu)=[\int f\,\mathrm{d}\mu][\int g\,\mathrm{d}\nu]$。
\end{enumerate}
\end{exercise}
% SOURCE p=84; printed=69; object=ch02-exercise-fubini-countable-factor; number=52; region=bottom
\begin{exercise}\label{exercise:ch02-fubini-countable-factor}
当 $(X,\mathscr{M},\mu)$ 为任意测度空间，$Y$ 为可数集，$\mathscr{N}=\mathscr{P}(Y)$，$\nu$ 为 $Y$ 上的计数测度时，Fubini--Tonelli 定理仍成立。（比较\cref{theorem:ch02-integral-series,theorem:ch02-absolutely-integrable-series}。）
\end{exercise}
\end{document}
